\chapter{Teori Ukuran}\label{ch:b3:measure}

Berapa panjang sebuah himpunan bagian $\R$? Jawaban naifnya --- yaitu memberi
setiap himpunan sebuah panjang yang invarian terhadap translasi dan memperluas panjang
selang --- ternyata \emph{mustahil}: konstruksi Vitali, pada
akhir bab ini, menghasilkan himpunan yang tak punya panjang yang
konsisten. Teori \omterm{def:b3:measure:measure}{ukuran} adalah kemunduran yang tertib: yakni membatasi
perhatian pada kelas kaya berisi himpunan \emph{terukur}, yang di atasnya
panjang yang aditif secara terbilang memang ada dan tunggal. Ganjarannya
luar biasa --- integral Lebesgue (\cref{ch:b3:lebesgue}),
ruang $L^p$ pada analisis fungsional, dan seluruh peluang
modern (\cref{ch:b3:probability}) dibangun di atas ketiga
teorema yang dibuktikan di sini: lema ketunggalan Dynkin,
teorema perluasan Carathéodory, dan keberadaan
\omterm{def:b3:measure:lebesgueouter}{ukuran Lebesgue}.

\section{Aljabar-$\sigma$}

\begin{definition}\label{def:b3:measure:sigmaalgebra}
Sebuah \emph{aljabar-$\sigma$}\index{aljabar-sigma@aljabar-$\sigma$}
pada himpunan $X$ adalah keluarga $\mathcal A$ berisi himpunan bagian yang memuat
$\varnothing$, stabil terhadap komplemen dan terhadap
gabungan \emph{terbilang} (sehingga juga terhadap irisan terbilang, selisih
himpunan, dan ia memuat $X$). Pasangan $(X, \mathcal A)$
disebut \emph{ruang terukur}; dan anggota $\mathcal A$ disebut
\emph{himpunan terukur}. Untuk sembarang keluarga $\mathcal E$ berisi himpunan bagian,
$\sigma(\mathcal E)$ melambangkan aljabar-$\sigma$ terkecil
yang memuat $\mathcal E$ (yakni irisan semuanya ---
sebab irisan aljabar-$\sigma$ tetap demikian).
\end{definition}

\begin{definition}\label{def:b3:measure:borel}
\emph{Aljabar-$\sigma$ Borel}\index{aljabar-sigma Borel@%
aljabar-$\sigma$ Borel} sebuah ruang \omterm{def:b3:topology:topology}{topologi} adalah
$\mathcal B(X) = \sigma(\{\text{himpunan terbuka}\})$. Pada $\R$:
$\mathcal B(\R)$ juga dibangkitkan oleh selang terbuka, oleh
selang tertutup, oleh sinar $(-\infty, a]$, dan oleh sinar
yang ujungnya rasional (\cref{exo:b3:measure:1}) --- sebab setiap
keluarga itu membangkitkan \omterm{def:b3:topology:topology}{himpunan terbukanya} lewat operasi terbilang, misalnya
setiap \omterm{def:b3:topology:topology}{himpunan terbuka} $\R$ berupa gabungan terbilang selang terbuka
yang datanya rasional.
\end{definition}

\begin{definition}\label{def:b3:measure:dynkin}
Sebuah \emph{sistem-$\pi$}\index{sistem-pi@sistem-$\pi$} adalah keluarga
yang stabil terhadap irisan berhingga. Sebuah
\emph{sistem-$\lambda$}\index{sistem-lambda@sistem-$\lambda$}
(kelas Dynkin) adalah keluarga $\mathcal D$ dengan: $X \in \mathcal
D$; $A, B \in \mathcal D$ dan $A \subseteq B$ $\Rightarrow$ $B
\setminus A \in \mathcal D$; serta $A_n \uparrow A$ dengan $A_n \in
\mathcal D$ $\Rightarrow$ $A \in \mathcal D$.
\end{definition}

\begin{theorem}[Lema $\pi$--$\lambda$ Dynkin]
\label{thm:b3:measure:dynkin}
Jika sebuah sistem-$\lambda$ $\mathcal D$ memuat sistem-$\pi$
$\mathcal P$, maka $\mathcal D \supseteq
\sigma(\mathcal P)$.\index{lema Dynkin}
\end{theorem}

\begin{proof}
Misalkan $\mathcal D_0$ sistem-$\lambda$ terkecil yang memuat
$\mathcal P$ (yakni irisan semuanya); cukup ditunjukkan bahwa
$\mathcal D_0$ merupakan aljabar-$\sigma$, sebab lalu
$\sigma(\mathcal P) \subseteq \mathcal D_0 \subseteq \mathcal
D$. Sistem-$\lambda$ yang stabil terhadap irisan berhingga
\emph{memang} aljabar-$\sigma$: sebab komplemennya ($X \setminus A = X
\setminus A$ dengan $A \subseteq X$), gabungan berhingganya ($A \cup B =
X \setminus ((X\setminus A)\cap(X\setminus B))$), dan gabungan terbilangnya
lewat $\bigcup_{k \leq n}A_k \uparrow \bigcup_kA_k$. Jadi kita
buktikan bahwa $\mathcal D_0$ merupakan sistem-$\pi$, dalam dua langkah. Misalkan
\[
\mathcal D_1 = \{A \in \mathcal D_0 : A \cap P \in \mathcal
D_0 \ \forall P \in \mathcal P\}.
\]
Keluarga $\mathcal D_1$ merupakan sistem-$\lambda$ (ketiga aksiomanya
diperiksa dengan mengiriskannya dengan $P$: misalnya $(B\setminus A)\cap P
= (B \cap P)\setminus(A \cap P)$, yaitu selisih sejati di dalam
$\mathcal D_0$) dan memuat $\mathcal P$ (karena sistem-$\pi$):
sehingga $\mathcal D_1 = \mathcal D_0$. Sekarang misalkan
\[
\mathcal D_2 = \{A \in \mathcal D_0 : A \cap D \in \mathcal
D_0\ \forall D \in \mathcal D_0\}.
\]
Menurut langkah sebelumnya, $\mathcal D_2 \supseteq \mathcal P$; dan
$\mathcal D_2$ merupakan sistem-$\lambda$ lewat pemeriksaan yang sama:
jadi $\mathcal D_2 = \mathcal D_0$, dan itu tepat mengatakan bahwa
$\mathcal D_0$ stabil terhadap irisan.
\end{proof}

\section{Ukuran}

\begin{definition}\label{def:b3:measure:measure}
Sebuah \emph{ukuran}\index{ukuran} pada $(X, \mathcal A)$ adalah pemetaan
$\mu \colon \mathcal A \to [0, +\infty]$ dengan $\mu(\varnothing)
= 0$ yang bersifat \emph{aditif-$\sigma$}: yakni untuk
$(A_n)_{n\in\N}$ yang saling lepas berpasangan,
\[
\mu\Bigl(\bigsqcup_n A_n\Bigr) = \sum_n \mu(A_n).
\]
Tripel $(X, \mathcal A, \mu)$ disebut \emph{ruang ukuran}; ukuran $\mu$ disebut
\emph{berhingga} jika $\mu(X) < \infty$, sebuah \emph{ukuran
peluang} jika $\mu(X) = 1$, dan \emph{berhingga-$\sigma$} jika $X$ merupakan
gabungan terbilang himpunan berukuran berhingga. Contohnya: ukuran
cacah pada $(\N, \mathcal P(\N))$; massa Dirac $\delta_a(A)
= \mathbf 1_{a \in A}$; dan, yang menjadi pokok bab ini,
\omterm{def:b3:measure:lebesgueouter}{ukuran Lebesgue}.
\end{definition}

\begin{proposition}\label{prop:b3:measure:basics}
Misalkan $\mu$ sebuah \omterm{def:b3:measure:measure}{ukuran}. (a) Kemonotonan: $A \subseteq B
\Rightarrow \mu(A) \leq \mu(B)$. (b) Subaditivitas terbilang:
$\mu(\bigcup A_n) \leq \sum\mu(A_n)$. (c) \omterm{def:b3:topology:continuity}{Kekontinuan} dari
bawah: $A_n \uparrow A \Rightarrow \mu(A_n) \to \mu(A)$. (d)
\omterm{def:b3:topology:continuity}{Kekontinuan} dari atas: $A_n \downarrow A$ \emph{dengan
$\mu(A_1) < \infty$} $\Rightarrow \mu(A_n) \to \mu(A)$.
\end{proposition}

\begin{proof}
(a) $B = A \sqcup (B\setminus A)$. (b) Saling lepaskan: $B_n = A_n
\setminus \bigcup_{k<n}A_k$ saling lepas dengan gabungan yang sama,
dan $\mu(B_n) \leq \mu(A_n)$. (c) $A = \bigsqcup_n (A_n
\setminus A_{n-1})$ (dengan $A_0 = \varnothing$): jumlah parsial
$\sum\mu(A_n\setminus A_{n-1})$ adalah $\mu(A_n)$. (d) Terapkan (c)
pada $A_1 \setminus A_n \uparrow A_1 \setminus A$ lalu kurangkan
dari $\mu(A_1)$ --- dan keberhinggaannya membuat pengurangannya
sah. Contoh penyangkal tanpa itu: $A_n = [n, \infty)$ untuk
\omterm{def:b3:measure:lebesgueouter}{ukuran Lebesgue}: $A_n \downarrow \varnothing$ padahal $\mu(A_n) =
\infty$.
\end{proof}

\begin{theorem}[Ketunggalan]\label{thm:b3:measure:uniqueness}
Misalkan $\mu, \nu$ \omterm{def:b3:measure:measure}{ukuran} pada $\sigma(\mathcal P)$, dengan $\mathcal
P$ sebuah sistem-$\pi$ dan $\mu = \nu$ pada $\mathcal P$. Jika ada
himpunan $P_k \in \mathcal P$ dengan $P_k \uparrow X$ dan
$\mu(P_k) < \infty$, maka $\mu = \nu$ pada seluruh
$\sigma(\mathcal P)$.
\end{theorem}

\begin{proof}
Tetapkan $k$ lalu tinjaulah \omterm{def:b3:measure:measure}{ukuran} berhingga $\mu_k(A) = \mu(A \cap
P_k)$ dan $\nu_k(A) = \nu(A \cap P_k)$ pada $\sigma(\mathcal
P)$: keduanya bersesuaian pada $\mathcal P$, sebab $P \cap P_k \in \mathcal
P$ (karena sistem-$\pi$), dan keduanya memberi $X$ nilai berhingga yang sama, yaitu
$\mu(P_k)$. Kelas $\mathcal D = \{A : \mu_k(A) =
\nu_k(A)\}$ merupakan sistem-$\lambda$: sebab $X \in \mathcal D$; selisih
sejatinya lewat pengurangan (karena nilainya berhingga); dan limit naiknya
lewat \omterm{def:b3:topology:continuity}{kekontinuan} dari bawah (\cref{prop:b3:measure:basics}(c)).
Ia memuat sistem-$\pi$ $\mathcal P$, sehingga Dynkin
(\cref{thm:b3:measure:dynkin}) memberikan $\mathcal D \supseteq
\sigma(\mathcal P)$: jadi $\mu_k = \nu_k$ di mana-mana. Akhirnya, untuk
sembarang $A \in \sigma(\mathcal P)$, \omterm{def:b3:topology:continuity}{kekontinuan} dari bawah sepanjang $A
\cap P_k \uparrow A$ memberikan $\mu(A) = \lim_k\mu_k(A) =
\lim_k\nu_k(A) = \nu(A)$.
\end{proof}

\section{Ukuran luar dan teorema Carathéodory}

\begin{definition}\label{def:b3:measure:outer}
Sebuah \emph{ukuran luar}\index{ukuran luar} pada $X$ adalah pemetaan
$\mu^* \colon \mathcal P(X) \to [0, \infty]$ dengan
$\mu^*(\varnothing) = 0$, yang monoton dan subaditif secara terbilang.
Sebuah himpunan $A$ disebut \emph{terukur-$\mu^*$} (menurut Carathéodory) jika
ia memecah setiap himpunan secara aditif:
\[
\mu^*(E) = \mu^*(E \cap A) + \mu^*(E \setminus A)
\qquad \text{untuk setiap } E \subseteq X
\]
($\leq$ selalu berlaku menurut subaditivitasnya; isinya adalah $\geq$).
\end{definition}

\begin{theorem}[Carathéodory]\label{thm:b3:measure:caratheodory}
Himpunan yang terukur-$\mu^*$ membentuk aljabar-$\sigma$ $\mathcal
M$, dan $\mu^*\restriction_{\mathcal M}$ merupakan \omterm{def:b3:measure:measure}{ukuran}.
Lebih lanjut setiap himpunan dengan $\mu^*(N) = 0$ termasuk $\mathcal M$
(sehingga \omterm{def:b3:measure:measure}{ukurannya} \emph{\omterm{def:b3:complete:complete}{lengkap}}).\index{teorema
Carathéodory}
\end{theorem}

\begin{proof}
Keluarga $\mathcal M$ memuat $\varnothing$ dan stabil terhadap
komplemen (sebab syarat pendefinisinya simetris terhadap $A$ dan
$X\setminus A$). \emph{Gabungan berhingga}: misalkan $A, B \in \mathcal
M$ dan $E$ sembarang; dengan memecah $E$ oleh $A$, lalu setiap kepingannya oleh
$B$:
\[
\mu^*(E) = \mu^*(E\cap A\cap B) + \mu^*(E\cap A\setminus B) +
\mu^*(E\cap B\setminus A) + \mu^*(E\setminus(A\cup B)).
\]
Ketiga kepingan pertamanya menyelimuti $E \cap (A \cup B)$, sehingga
subaditivitasnya memberikan $\mu^*(E) \geq \mu^*(E\cap(A\cup B)) +
\mu^*(E\setminus(A\cup B))$: jadi $A \cup B \in \mathcal M$. Lewat
induksi, gabungan berhingganya pun demikian; dan bersama komplemennya, manipulasi
kesalinglepasan yang berhingga menjadi tersedia.

\emph{Keaditifan pada $\mathcal M$}: untuk $A, B \in
\mathcal M$ yang saling lepas dan $E$ apa pun: $\mu^*(E\cap(A\sqcup B)) = \mu^*(E\cap
A) + \mu^*(E \cap B)$ (dengan memecahnya oleh $A$); lalu lewat induksi,
\[
\mu^*\Bigl(E \cap \bigsqcup_{k\leq n}A_k\Bigr)
= \sum_{k\leq n}\mu^*(E\cap A_k).
\tag{$*$}
\]

\emph{Gabungan terbilang}: misalkan $(A_k) \subseteq \mathcal M$
saling lepas (dan ini sudah cukup, lewat pelepasan di dalam aljabar
$\mathcal M$), $A = \bigsqcup A_k$, serta $E$ sembarang. Dengan memakai
$\bigsqcup_{k\leq n}A_k \in \mathcal M$ dan kemonotonannya:
\[
\mu^*(E) = \mu^*\Bigl(E\cap\bigsqcup_{k\leq n}A_k\Bigr) +
\mu^*\Bigl(E\setminus\bigsqcup_{k\leq n}A_k\Bigr)
\geq \sum_{k \leq n}\mu^*(E\cap A_k) + \mu^*(E\setminus A)
\]
menurut ($*$). Biarkan $n \to \infty$ lalu pakai subaditivitas terbilangnya
secara terbalik:
\[
\mu^*(E) \geq \sum_{k}\mu^*(E\cap A_k) + \mu^*(E\setminus A)
\geq \mu^*(E \cap A) + \mu^*(E\setminus A) \geq \mu^*(E):
\]
sehingga semua ketaksamaannya menjadi kesamaan. Ini membuktikan sekaligus $A \in
\mathcal M$ dan, dengan mengambil $E = A$, keaditifan terbilang
$\mu^*$ pada $\mathcal M$.

\emph{Himpunan nol}: jika $\mu^*(N) = 0$, maka untuk $E$ apa pun:
$\mu^*(E \cap N) + \mu^*(E\setminus N) \leq 0 + \mu^*(E)$:
jadi $N \in \mathcal M$.
\end{proof}

\section{Ukuran Lebesgue pada \texorpdfstring{$\R$}{R}}

\begin{definition}\label{def:b3:measure:lebesgueouter}
\emph{Ukuran luar Lebesgue}\index{ukuran Lebesgue} dari $A
\subseteq \R$ adalah
\[
\lambda^*(A) = \inf\Bigl\{\sum_{n} (b_n - a_n) :
A \subseteq \bigcup_n \intoo{a_n}{b_n}\Bigr\}
\]
(dengan selimut terbilang oleh selang terbuka).
\end{definition}

\begin{lemma}\label{lem:b3:measure:outerbasics}
Fungsi $\lambda^*$ merupakan \omterm{def:b3:measure:outer}{ukuran luar} yang invarian terhadap translasi,
dan $\lambda^*(I) = \ell(I)$ (yakni panjangnya) untuk setiap selang
$I$.
\end{lemma}

\begin{proof}
\omterm{def:b3:measure:outer}{Ukuran luarnya}: $\varnothing$ terselimuti oleh selang yang sekecil
apa pun; kemonotonannya jelas; dan subaditivitasnya: diberikan selimut
setiap $A_n$ dalam $\varepsilon 2^{-n}$ dari infimumnya,
gabungannya menyelimuti $\bigcup A_n$ dengan panjang total $\leq \sum
\lambda^*(A_n) + \varepsilon$. Keinvarianan translasinya:
translasikan selimutnya.

Untuk panjangnya: cukup ditangani $I = \intcc ab$ (sebab jenis lainnya
berbeda pada ujungnya, yang berukuran luar $0$: selimuti dengan
selang mungil; lalu himpitlah lewat perbandingan bertipe $\intcc{a+\varepsilon}{b -
\varepsilon} \subseteq \intoo ab$).
Untuk $\lambda^*(\intcc ab) \leq b - a$: selimuti dengan
$\intoo{a-\varepsilon}{b+\varepsilon}$. Sebaliknya misalkan $\intcc
ab \subseteq \bigcup_n\intoo{a_n}{b_n}$: menurut
\emph{kekompakannya} (Borel--Lebesgue,
\cref{thm:b3:topology:metriccompact}), berhingga selang
sudah cukup, sebutlah $I_1, \dots, I_N$. Kita tunjukkan $\sum_{k\leq N}(b_k -
a_k) \geq b - a$ lewat induksi pada $N$: pilihlah $I_{k_1} \ni a$; jika
$b_{k_1} > b$ maka selesai (sebab $b_{k_1} - a_{k_1} > b - a$); jika tidak,
ruas $\intcc{b_{k_1}}b$ terselimuti oleh $N - 1$ selang
sisanya, dan induksinya memberikan $\sum_{k \neq k_1}(b_k - a_k)
\geq b - b_{k_1}$, sedangkan $b_{k_1} - a_{k_1} > b_{k_1} - a$:
lalu jumlahkan.
\end{proof}

\begin{theorem}[Ukuran Lebesgue]\label{thm:b3:measure:lebesgue}
Setiap himpunan Borel di $\R$ bersifat terukur-$\lambda^*$.
Pembatasan $\lambda$ dari $\lambda^*$ pada aljabar-$\sigma$
$\mathcal L = \mathcal M_{\lambda^*} \supseteq \mathcal B(\R)$
(yakni \emph{aljabar-$\sigma$ Lebesgue}) merupakan satu-satunya \omterm{def:b3:measure:measure}{ukuran} pada
$\mathcal B(\R)$ yang memberi setiap selang panjangnya; ia
invarian terhadap translasi dan berhingga-$\sigma$.
\end{theorem}

\begin{proof}
Menurut \cref{thm:b3:measure:caratheodory} cukup ditunjukkan bahwa setiap
sinar $A = \intoo{-\infty}c$ bersifat terukur-$\lambda^*$
(sebab sinar membangkitkan $\mathcal B$,
\cref{def:b3:measure:borel}). Misalkan $E \subseteq \R$ dengan
$\lambda^*(E) < \infty$ dan $\bigcup I_n \supseteq E$ sebuah selimut
dengan $\sum\ell(I_n) \leq \lambda^*(E) + \varepsilon$. Setiap
$I_n$ terpecah menjadi dua selang $I_n' = I_n \cap A$ dan
$I_n'' = I_n\setminus A$ (sebab selang dikurangi sinar tetap
selang) dengan $\ell(I_n') + \ell(I_n'') = \ell(I_n)$; lalu
$I_n'$ menyelimuti $E \cap A$ dan $I_n''$ menyelimuti $E \setminus A$
(perbesarlah masing-masing menjadi selang terbuka berpanjang $\ell +
\varepsilon2^{-n}$ agar tetap sesuai definisinya), sehingga
\[
\lambda^*(E\cap A) + \lambda^*(E\setminus A)
\leq \sum_n\bigl(\ell(I_n') + \ell(I_n'')\bigr) + 2\varepsilon
\leq \lambda^*(E) + 3\varepsilon .
\]
Ketunggalannya: dua \omterm{def:b3:measure:measure}{ukuran} yang bersesuaian dengan panjang pada
sistem-$\pi$ berisi selang $\intoc ab$ (dan berhingga di sana) bersesuaian pada
$\sigma(\text{selang}) = \mathcal B$ menurut
\cref{thm:b3:measure:uniqueness} dengan $P_k = \intoc{-k}k$.
Keberhinggaan-$\sigma$-nya: $\R = \bigcup(-k, k]$.
\end{proof}

\begin{theorem}[Keteraturan]\label{thm:b3:measure:regularity}
Untuk setiap $A \in \mathcal L$:\index{keteraturan ukuran
Lebesgue}
\[
\lambda(A) = \inf\{\lambda(U) : U \supseteq A \text{ terbuka}\}
= \sup\{\lambda(K) : K \subseteq A \text{ kompak}\}.
\]
\end{theorem}

\begin{proof}
\emph{Sisi luarnya}: selimut $\bigcup I_n$ dengan $\sum\ell(I_n) \leq
\lambda(A) + \varepsilon$ merupakan \omterm{def:b3:topology:topology}{himpunan terbuka} $U \supseteq A$ dengan
$\lambda(U) \leq \lambda(A) + \varepsilon$ (lewat subaditivitasnya); dan jika
$\lambda(A) = \infty$ pernyataannya trivial. \emph{Sisi dalamnya}:
lebih dahulu misalkan $A$ terbatas, $A \subseteq [-M, M]$. Pilihlah sebuah
$U \supseteq ([-M,M]\setminus A)$ yang terbuka dengan $\lambda(U) \leq
\lambda([-M,M]\setminus A) + \varepsilon$; maka $K = [-M,
M]\setminus U$ bersifat \omterm{def:b3:topology:compact}{kompak}, $K \subseteq A$, dan
\[
\lambda(K) = \lambda([-M,M]) - \lambda([-M,M]\cap U)
\geq \lambda([-M,M]) - \bigl(\lambda([-M,M]) - \lambda(A) +
\varepsilon\bigr) = \lambda(A) - \varepsilon .
\]
Untuk $A$ yang umum: $\lambda(A) = \lim_M\lambda(A\cap[-M,M])$
(lewat \omterm{def:b3:topology:continuity}{kekontinuan} dari bawah) lalu terapkan kasus terbatasnya di dalamnya.
\end{proof}

\begin{example}\label{ex:b3:measure:cantor}
Himpunan Cantor (\cref{exo:b3:topology:10}) mempunyai $\lambda(C) = 0$:
sebab $C \subseteq C_n$, yaitu gabungan $2^n$ selang berpanjang
$3^{-n}$, sehingga $\lambda(C) \leq (2/3)^n \to 0$. Jadi ada himpunan nol
yang tak terbilang --- kardinalitas tak melihat \omterm{def:b3:measure:measure}{ukuran}. Sebaliknya,
\emph{himpunan Cantor gemuk} (\cref{exo:b3:measure:5}) bersifat tak padat
di mana pun tetapi berukuran positif: jadi \omterm{def:b3:topology:topology}{topologi} pun tak melihat \omterm{def:b3:measure:measure}{ukuran}.
Soal akhir pekan mendorong permainan ini sampai kesimpulannya yang
mencolok: bahwa ada himpunan terukur Lebesgue yang
bukan himpunan Borel.
\end{example}

\begin{theorem}[Vitali]\label{thm:b3:measure:vitali}
Tak ada \omterm{def:b3:measure:measure}{ukuran} pada \emph{seluruh} himpunan bagian $\R$ yang
invarian terhadap translasi dan memberi setiap selang panjangnya.
Khususnya $\mathcal L \neq \mathcal P(\R)$: jadi himpunan tak terukur
memang ada.\index{himpunan tak terukur Vitali}
\end{theorem}

\begin{proof}
Andaikan $\mu$ salah satunya. Pada $\intcc01$, tinjaulah relasi ekuivalensi
$x \sim y \iff x - y \in \Q$; lalu lewat \emph{aksioma pilihan}
ambillah satu wakil di $\intcc01$ per kelasnya: sebutlah himpunan $V$. Untuk
$q \in \Q\cap\intcc{-1}1$, translasi $V + q$ saling lepas
berpasangan (sebab dua titik $V$ yang berselisih rasional akan
ekuivalen padahal keduanya wakil yang berbeda) dan
\[
\intcc01 \subseteq \bigsqcup_{q \in \Q\cap\intcc{-1}1}(V + q)
\subseteq \intcc{-1}2 :
\]
dengan inklusi pertamanya karena setiap $x \in \intcc01$ berselisih dari
wakilnya $v$ sebesar bilangan rasional $q = x - v \in
\intcc{-1}1$. Lalu kemonotonan dan keaditifan-$\sigma$-nya memberikan
\[
1 \leq \sum_{q}\mu(V + q) \leq 3,
\qquad\text{dengan } \mu(V + q) = \mu(V) \text{ untuk setiap } q .
\]
Jumlah tak berhingga atas konstanta $\mu(V)$ bernilai $0$ atau $\infty$:
sehingga kedua batasnya tak mungkin berlaku bersamaan. Jadi $\mu$ semacam itu tak ada --- dan $V
\notin \mathcal L$, sebab $\lambda$ pada $\mathcal L$ mempunyai semua
sifat yang dipakai tadi.
\end{proof}

\begin{remark}\label{rem:b3:measure:banachtarski}
Di $\R^3$ kegagalannya lebih dramatis: paradoks Banach--Tarski
menguraikan sebuah bola menjadi lima kepingan yang dirakit kembali, lewat
rotasi dan translasi, menjadi \emph{dua} bola berjari-jari
sama --- sehingga bahkan volume yang aditif berhingga dan invarian terhadap rotasi
pada seluruh himpunan bagian $\R^3$ pun tak ada. Kepingannya, tentu saja,
tidak terukur. Jadi keterukuran bukanlah kehati-hatian
birokratis; ia adalah batas kekoherenan.
\end{remark}

\begin{method}\label{met:b3:measure:goodsets}
\emph{Asas himpunan baik}: untuk membuktikan bahwa semua himpunan
$\sigma(\mathcal E)$ mempunyai sebuah sifat, tunjukkan bahwa himpunan baiknya
membentuk aljabar-$\sigma$ (atau sistem-$\lambda$, bila sifatnya
bersifat teori \omterm{def:b3:measure:measure}{ukuran} dan $\mathcal E$ merupakan sistem-$\pi$ ---
lalu pakai Dynkin) yang memuat $\mathcal E$. Hampir setiap bukti pada
bab ini dan bab berikutnya merupakan contohnya. Untuk membuktikan dua \omterm{def:b3:measure:measure}{ukuran}
sama: periksalah keduanya pada sistem-$\pi$ pembangkitnya ditambah
keberhinggaan-$\sigma$ (\cref{thm:b3:measure:uniqueness}). Untuk
membangun sebuah \omterm{def:b3:measure:measure}{ukuran}: bangunlah \omterm{def:b3:measure:outer}{ukuran luar} lewat selimut lalu kutiplah
Carathéodory.
\end{method}

\section{Latihan}

\begin{exercise}[$\star$]\label{exo:b3:measure:1}
(a) Tunjukkan bahwa $\{A \subseteq X : A$ atau $X\setminus A$ bersifat
terbilang$\}$ merupakan aljabar-$\sigma$: yakni yang dibangkitkan oleh
himpunan tunggal.
(b) Tunjukkan bahwa $\mathcal B(\R)$ dibangkitkan oleh masing-masing berikut: selang
terbuka; selang tertutup; sinar $\intoc{-\infty}a$; dan sinar
dengan $a \in \Q$.
(c) Apakah keluarga berisi \emph{gabungan berhingga selang yang saling lepas}
$\intoc ab$ merupakan aljabar-$\sigma$? Apakah ia aljabar (yakni stabil terhadap
komplemen dan gabungan berhingga)?
\end{exercise}

\begin{exercise}[$\star$]\label{exo:b3:measure:2}
(a) Buktikan inklusi--eksklusi bagi \omterm{def:b3:measure:measure}{ukuran} berhingga: $\mu(A\cup
B) = \mu(A) + \mu(B) - \mu(A\cap B)$, beserta versi
tiga himpunannya.
(b) Berilah contoh yang menunjukkan bahwa \omterm{def:b3:topology:continuity}{kekontinuan} dari atas
(\cref{prop:b3:measure:basics}(d)) gagal tanpa
pengandaian keberhinggaannya.
(c) Tunjukkan bahwa himpunan terbilang berukuran Lebesgue nol.
Lalu simpulkan $\lambda(\Q) = 0$ dan $\lambda(\intcc01\setminus\Q) =
1$.
\end{exercise}

\begin{exercise}[$\star\star$]\label{exo:b3:measure:3}
Misalkan $\mu, \nu$ \omterm{def:b3:measure:measure}{ukuran} peluang pada $\mathcal B(\R)$
dengan $\mu(\intoc{-\infty}t) = \nu(\intoc{-\infty}t)$ untuk setiap
$t \in \R$. Tunjukkan $\mu = \nu$. (Hal ini menjadikan
\emph{fungsi distribusi} $F(t) = \mu(\intoc{-\infty}t)$ sebuah
invarian yang \omterm{def:b3:complete:complete}{lengkap} --- yakni landasan
\cref{ch:b3:probability}.)
\end{exercise}

\begin{exercise}[$\star\star$]\label{exo:b3:measure:4}
(Borel--Cantelli, versi \omterm{def:b3:measure:measure}{ukuran}) Misalkan $(A_n)$ terukur
dengan $\sum_n\mu(A_n) < \infty$, dan $\limsup A_n =
\bigcap_N\bigcup_{n\geq N}A_n$ (yakni titik yang termasuk
tak berhingga banyak $A_n$). Tunjukkan $\mu(\limsup A_n) = 0$.
Terapannya: untuk hampir setiap $x \in \intcc01$, hanya berhingga
banyak $n$ yang memenuhi $\abs{x - p/q_n} \leq 4^{-n}$ bagi
bilangan rasional ke-$n$, yakni $p/q_n$, pada suatu pencacahan
$\Q\cap\intcc01$.
\end{exercise}

\begin{exercise}[$\star\star$]\label{exo:b3:measure:5}
(Himpunan Cantor gemuk) Ulangi konstruksi Cantor pada $\intcc01$,
tetapi pada langkah $n$ buanglah dari masing-masing $2^{n-1}$ selangnya sebuah
selang terbuka \emph{terpusat} berpanjang $4^{-n}$ saja. Tunjukkan
bahwa $K = \bigcap K_n$ yang dihasilkan bersifat \omterm{def:b3:topology:compact}{kompak}, berinterior
kosong (sebab tak ada selang yang bertahan), dan
\[
\lambda(K) = 1 - \sum_{n\geq1}2^{n-1}4^{-n} = \tfrac12 :
\]
yakni himpunan tak padat di mana pun yang berukuran $\frac12$. Lalu simpulkan adanya
himpunan bagian $\intcc01$ yang \emph{\omterm{rem:b3:complete:meagre}{kurus}} tetapi berukuran penuh $1$, dan sebuah
himpunan bagian terbuka yang padat dan berukuran $< \varepsilon$.
\end{exercise}

\begin{exercise}[$\star\star$]\label{exo:b3:measure:6}
Misalkan $\mu$ sebuah \omterm{def:b3:measure:measure}{ukuran} pada $\mathcal B(\R)$ yang invarian terhadap
translasi, dengan $c = \mu(\intoc01) < \infty$. Tunjukkan $\mu =
c\,\lambda$ pada $\mathcal B(\R)$. \emph{(Hitunglah $\mu$ pada selang diadik dengan membagi
$\intoc01$ menjadi $2^n$ translasinya, lalu panggillah
\cref{thm:b3:measure:uniqueness}.)}
\end{exercise}

\begin{exercise}[$\star\star$]\label{exo:b3:measure:7}
(Penghampiran) Misalkan $A \in \mathcal L$ dengan $\lambda(A) <
\infty$ dan $\varepsilon > 0$. Tunjukkan bahwa ada gabungan
\emph{berhingga} selang $B$ dengan $\lambda(A\,\triangle\,B) <
\varepsilon$ (dengan $\triangle$ = selisih simetris).
\emph{(Lewat keteraturannya: himpitlah $K \subseteq A \subseteq U$ lalu pakai
struktur $U$ yang terbuka sebagai gabungan terbilang
selang, ditambah kekompakan $K$.)}
\end{exercise}

\begin{exercise}[$\star\star\star$]\label{exo:b3:measure:8}
(Steinhaus) Misalkan $A \in \mathcal L$ dengan $\lambda(A) > 0$. Tunjukkan
bahwa $A - A = \{x - y : x, y \in A\}$ memuat sebuah selang
di sekitar $0$.
\emph{(Reduksikan ke $\lambda(A) < \infty$; lalu lewat keteraturan
bergaya \cref{exo:b3:measure:7}, carilah selang $I$
dengan $\lambda(A \cap I) > \frac34\ell(I)$; maka untuk $\abs t <
\frac12\ell(I)$, himpunan $A\cap I$ dan $(A\cap I) + t$ keduanya
duduk di selang berpanjang $\frac32\ell(I)$ dan berukuran total
$> \frac32\ell(I)$: sehingga keduanya harus berpotongan.)}
\end{exercise}

\begin{exercise}[$\star\star\star$]\label{exo:b3:measure:9}
Tunjukkan bahwa setiap $A \in \mathcal L$ dengan $\lambda(A) > 0$
memuat himpunan bagian yang tak terukur. \emph{(Iriskan $A$ dengan
translasi $V + q$ dari himpunan Vitali: sebab jika semua $A \cap (V+q)$
terukur, masing-masing akan berukuran nol menurut argumen
\cref{thm:b3:measure:vitali} --- dan Steinhaus
(\cref{exo:b3:measure:8}) menolong: himpunan terukur berukuran positif
di dalam $V + q$ akan membuat $(V+q) - (V+q) \supseteq$
sebuah selang, dan itu bertentangan dengan kenyataan bahwa himpunan selisih itu memotong $\Q$
hanya di $0$; lalu simpulkan lewat subaditivitasnya.)}
\end{exercise}

\begin{exercise}[$\star\star$]\label{exo:b3:measure:10}
Tunjukkan bahwa $A \subseteq \R$ dengan $\lambda^*(A) < \infty$ bersifat
terukur Lebesgue jika dan hanya jika untuk setiap $\varepsilon > 0$ ada
$U \supseteq A$ yang terbuka dengan $\lambda^*(U \setminus A) <
\varepsilon$, dan jika dan hanya jika ada himpunan $G_\delta$ $G \supseteq A$
dengan $\lambda^*(G\setminus A) = 0$. (Jadi himpunan Lebesgue merupakan
himpunan Borel modulo himpunan nol.)
\end{exercise}

\begin{exercise}[$\star\star$]\label{exo:b3:measure:11}
(\omterm{def:b3:topology:continuity}{Kekontinuan} sepanjang limit monoton, dan ketajamannya)
(a) Tunjukkan bahwa untuk himpunan terukur berlaku $\mu(\liminf A_n) \leq
\liminf\mu(A_n)$ (yakni Fatou untuk himpunan), dan bahwa jika
$\mu\bigl(\bigcup A_n\bigr) < \infty$, maka juga
$\limsup\mu(A_n) \leq \mu(\limsup A_n)$.
(b) Tampilkan, untuk \omterm{def:b3:measure:lebesgueouter}{ukuran Lebesgue} pada $\R$, sebuah barisan dengan
$\mu(A_n) = 1$ untuk setiap $n$ namun $\mu(\limsup A_n) = 0$: jadi
hipotesis keberhinggaan pada ketaksamaan kedua bukanlah hiasan.
(c) Simpulkan: jika $\sum\mu(A_n) < \infty$ maka
$\mu(\limsup A_n) = 0$ (Borel--Cantelli lagi), dan jika
$A_n$ naik atau turun (dengan $\mu(A_1) < \infty$ pada kasus
turunnya), maka $\mu(\lim A_n) = \lim\mu(A_n)$.
\end{exercise}

\begin{exercise}[$\star\star\star$]\label{exo:b3:measure:12}
(Teorema Egorov) Misalkan $\mu(X) < \infty$ dan $f_n \to f$
titik demi titik, dengan semuanya terukur (bernilai real). Untuk $k, N \geq 1$
tetapkan
\[
E_{k,N} = \bigcap_{n \geq N}\Bigl\{x : \abs{f_n(x) - f(x)}
\leq \tfrac1k\Bigr\} .
\]
(a) Tunjukkan bahwa untuk $k$ yang tetap, $E_{k,N} \nearrow X$ ketika $N \to
\infty$, lalu simpulkan adanya $N_k$ dengan $\mu(X \setminus E_{k,N_k})
\leq \varepsilon2^{-k}$.
(b) Simpulkan \emph{teorema Egorov}: bahwa untuk setiap $\varepsilon
> 0$ ada $A$ yang terukur dengan $\mu(X\setminus A) \leq
\varepsilon$ sedemikian sehingga $f_n \to f$ \emph{secara seragam pada $A$}
--- jadi kekonvergenan titik demi titik menjadi kekonvergenan seragam di luar
himpunan yang sekecil apa pun.
(c) Tunjukkan bahwa teoremanya gagal pada $(\R, \lambda)$: gundukan
berjalan $f_n = \mathbf 1_{\intcc n{n+1}}$ konvergen titik demi titik
ke $0$ tetapi tidak seragam pada komplemen himpunan berukuran
berhingga mana pun. Di manakah (a) memakai $\mu(X) < \infty$?
\end{exercise}

\section{Soal: tangga Cantor--Vitali dan sebuah himpunan
terukur yang bukan Borel}

\begin{omfigure}
\begin{tikzpicture}[scale=5.2]
  \draw[->] (0,0) -- (1.06,0) node[right, font=\small] {$x$};
  \draw[->] (0,0) -- (0,1.06);
  \draw[black!30] (1,0) -- (1,1) -- (0,1);
  \draw[omThm, thick]
    (0,0)
    -- (0.037,0.125) -- (0.074,0.125)
    -- (0.111,0.25)  -- (0.222,0.25)
    -- (0.259,0.375) -- (0.296,0.375)
    -- (0.333,0.5)   -- (0.667,0.5)
    -- (0.704,0.625) -- (0.741,0.625)
    -- (0.778,0.75)  -- (0.889,0.75)
    -- (0.926,0.875) -- (0.963,0.875)
    -- (1,1);
  \node[font=\small, omThm] at (0.32,0.86)
    {$c$ (hampiran kedalaman 3)};
\end{tikzpicture}

{\small Tangga Cantor--Vitali: yang konstan pada setiap celah
himpunan Cantor, namun mendaki dari $0$ ke $1$ secara \omterm{def:b3:topology:continuity}{kontinu}.
Turunannya lenyap hampir di mana-mana --- jadi seluruh pendakiannya
terjadi pada sebuah himpunan nol.}
\end{omfigure}

\begin{problem}[{Soal akhir pekan --- tangga setan,
dan $\mathcal B(\R) \subsetneq \mathcal L$}]
\label{pb:b3:measure:1}
Kita bangun \emph{fungsi Cantor--Vitali} (yakni tangga
setan), memakainya untuk mengangkut \omterm{def:b3:measure:measure}{ukuran} secara patologis, lalu
menutup dengan teorema yang tak diberikan oleh argumen lunak mana pun: bahwa
ada himpunan terukur Lebesgue yang bukan Borel. Notasinya:
$C$ adalah himpunan Cantor, $C_n$ tahap ke-$n$-nya ($2^n$ selang
berpanjang $3^{-n}$), dan setiap $x \in C$ berangka terner
$x = \sum 2b_n3^{-n}$ dengan $b_n \in \{0,1\}$
(\cref{exo:b3:topology:10}).

\medskip
\noindent\textbf{Bagian I --- Tangganya.} Definisikan $c_0(x) = x$
lalu $c_{n+1}$ dari $c_n$ lewat
\[
c_{n+1}(x) = \begin{cases}
\tfrac12\,c_n(3x) & 0 \leq x \leq \tfrac13,\\[2pt]
\tfrac12 & \tfrac13 \leq x \leq \tfrac23,\\[2pt]
\tfrac12 + \tfrac12\,c_n(3x - 2) & \tfrac23 \leq x \leq 1.
\end{cases}
\]
\begin{enumerate}
  \item Tunjukkan bahwa setiap $c_n$ \omterm{def:b3:topology:continuity}{kontinu}, tak turun, dengan
        $c_n(0) = 0$, $c_n(1) = 1$, dan bahwa
        $\norm{c_{n+1} - c_n}_\infty \leq
        \tfrac12\norm{c_n - c_{n-1}}_\infty$.
  \item Simpulkan bahwa $(c_n)$ konvergen seragam ke sebuah $c$ yang
        \omterm{def:b3:topology:continuity}{kontinu} dan tak turun dengan $c(0) = 0$, $c(1) = 1$
        (yakni \emph{fungsi} Cantor--Vitali), yang memenuhi
        relasi serupa-diri yang sama seperti $c_{n+1}$ di atas.
  \item Tunjukkan bahwa $c$ konstan pada setiap \omterm{def:b3:topology:components}{komponen terhubung}
        $\intcc01\setminus C$, dan bahwa untuk $x = \sum_n
        2b_n3^{-n} \in C$ berlaku $c(x) = \sum_n b_n2^{-n}$: jadi
        tangganya membaca angka Cantor dalam biner
        (yakni fungsi $g$ pada \cref{pb:b3:topology:1}, yang dibuat monoton
        dan global).
  \item Simpulkan bahwa $c$ dapat diturunkan, dengan $c' = 0$, di
        setiap titik $\intcc01\setminus C$: sehingga $c' = 0$
        \emph{hampir di mana-mana terhadap $\lambda$}
        (\cref{ex:b3:measure:cantor}). Lalu simpulkan bahwa
        teorema fundamental kalkulus gagal bagi $c$:
        \[
        c(1) - c(0) = 1 \neq 0 = \int_0^1 c'(t)\,\dd t
        \]
        (dengan integralnya atas himpunan berukuran penuh tempat
        $c' = 0$; dan dengan mendahului \cref{ch:b3:lebesgue}, himpunan nol
        tak memengaruhi integral). Hipotesis mana pada
        teorema fundamental $\mathcal C^1$ yang dilanggar?
  \item Tunjukkan bahwa $c(C) = \intcc01$: jadi himpunan nol $C$ terpetakan
        \emph{pada} himpunan berukuran penuh.
\end{enumerate}

\medskip
\noindent\textbf{Bagian II --- \omterm{def:b3:topology:continuity}{Homeomorfisma} yang bengkok.} Misalkan
$h(x) = \frac{c(x) + x}{2}$.
\begin{enumerate}[resume]
  \item Tunjukkan bahwa $h \colon \intcc01 \to \intcc01$ merupakan
        \omterm{def:b3:topology:continuity}{homeomorfisma} (yakni naik sejati, \omterm{def:b3:topology:continuity}{kontinu}, dan
        surjektif).
  \item Tunjukkan bahwa $\lambda\bigl(h(\intcc01\setminus C)\bigr) =
        \tfrac12$: sebab pada setiap celah berpanjang $\ell$, $h$ bekerja sebagai
        pemetaan afin berkemiringan $\tfrac12$, dan celahnya berpanjang
        total $1$.
  \item Simpulkan $\lambda\bigl(h(C)\bigr) = \tfrac12$: jadi
        peta \omterm{def:b3:topology:continuity}{homeomorfik} sebuah himpunan nol dapat berukuran positif.
        (Di manakah hal ini bertentangan dengan gerak hati yang naif
        tentang ``\omterm{def:b3:measure:measure}{ukuran}''?)
\end{enumerate}

\medskip
\noindent\textbf{Bagian III --- Sebuah himpunan terukur yang bukan
Borel.}
\begin{enumerate}[resume]
  \item Menurut \cref{exo:b3:measure:9}, pilihlah $W
        \subseteq h(C)$ yang tak terukur. Tunjukkan bahwa $Z = h^{-1}(W) \subseteq C$
        bersifat terukur Lebesgue. \emph{(Sebab ia himpunan bagian sebuah himpunan
        nol; lihat \omterm{def:b3:complete:complete}{kelengkapannya}, \cref{thm:b3:measure:caratheodory}.)}
  \item Tunjukkan bahwa prapeta himpunan Borel di bawah \omterm{def:b3:topology:continuity}{pemetaan
        kontinu} bersifat Borel. \emph{(Lewat asas himpunan baik:
        $\{B : h^{-1}(B) \in \mathcal B\}$ merupakan
        aljabar-$\sigma$ yang memuat \omterm{def:b3:topology:topology}{himpunan terbukanya} --- perhatikan
        arah pemetaannya.)}
  \item Simpulkan bahwa $Z$ \emph{bukan} Borel: sebab seandainya demikian,
        $W = (h^{-1})^{-1}(Z)$ akan menjadi Borel (terapkan pertanyaan
        10 pada $h^{-1}$ yang \omterm{def:b3:topology:continuity}{kontinu}), sehingga terukur ---
        kontradiksi. Karena itu
        \[
        \boxed{\ \mathcal B(\R) \subsetneq \mathcal L\ }
        \]
        dan \omterm{def:b3:complete:complete}{kelengkapan} sungguh memperbesar dunia Borel.
  \item Tampilkan sebuah fungsi $g$ yang terukur Lebesgue dan sebuah
        fungsi \omterm{def:b3:topology:continuity}{kontinu} $\varphi$ sedemikian sehingga $g \circ
        \varphi$ tak terukur Lebesgue: jadi keterukuran,
        tak seperti \omterm{def:b3:topology:continuity}{kekontinuan}, tak terkomposisikan. \emph{(Ambil $g =
        \mathbf 1_Z$ dan $\varphi = h^{-1}$, sambil mendahului
        definisi fungsi terukur pada
        \cref{ch:b3:lebesgue}: bahwa prapeta himpunan Borel merupakan
        himpunan Lebesgue. Di manakah kita harus berhati-hati tentang
        aljabar-$\sigma$ mana yang dipakai pada sasarannya?)}
\end{enumerate}

\medskip
\noindent\textbf{Bagian IV --- \omterm{def:b3:topology:interior}{Penutup}.}
\begin{enumerate}[resume]
  \item Urutkan kelas berikut menurut inklusi sejatinya lalu
        benarkan setiap kesejatiannya dengan contoh dari bab
        ini beserta soalnya: himpunan terbilang; himpunan nol
        Borel; himpunan nol Lebesgue; himpunan Borel; himpunan Lebesgue;
        dan himpunan sembarang.
\end{enumerate}

\medskip
\noindent\textbf{Bagian V --- \omterm{def:b3:measure:measure}{Ukuran} Cantor: massa pada sebuah himpunan
nol.} Tangga itu merupakan fungsi distribusi sebuah \omterm{def:b3:measure:measure}{ukuran} yang
luar biasa, yang kini kita bangun dengan perkakas bab ini
sendiri.
\begin{enumerate}[resume]
  \item (Lebesgue--Stieltjes, keberadaannya) Misalkan $F\colon\R\to\R$
        tak turun, \omterm{def:b3:topology:continuity}{kontinu}, dan terbatas. Pada selang setengah
        terbuka definisikan $\rho\bigl(\intoc ab\bigr) = F(b) -
        F(a)$ lalu, untuk $A \subseteq \R$,
        \[
        \mu_F^*(A) = \inf\Bigl\{\sum_k\bigl(F(b_k) -
        F(a_k)\bigr) : A \subseteq
        \bigcup_k\intoc{a_k}{b_k}\Bigr\} .
        \]
        Tunjukkan bahwa $\mu_F^*$ merupakan \omterm{def:b3:measure:outer}{ukuran luar} dan bahwa
        $\mu_F^*\bigl(\intoc ab\bigr) = F(b) - F(a)$
        \emph{(tirulah argumen kekompakan pada
        \cref{thm:b3:measure:lebesgue}, dengan memperbesar setiap
        $\intoc{a_k}{b_k}$ menjadi selang terbuka berbiaya-$F$
        $\leq \varepsilon2^{-k}$ --- di manakah
        \omterm{def:b3:topology:continuity}{kekontinuan} $F$ dipakai?)}.
  \item Tunjukkan bahwa setiap himpunan Borel bersifat
        terukur-$\mu_F^*$ dalam arti Carathéodory
        \emph{(seperti pada kasus Lebesgue, cukup diuji
        pada setengah garis; ikutilah bukti
        penerapan \cref{thm:b3:measure:caratheodory})}, sehingga
        $\mu_F = \mu_F^*$ yang dibatasi pada $\mathcal B(\R)$
        merupakan \omterm{def:b3:measure:measure}{ukuran} dengan $\mu_F(\intoc ab) = F(b) - F(a)$:
        yakni \emph{\omterm{def:b3:measure:measure}{ukuran} Lebesgue--Stieltjes} dari $F$.
  \item Terapkan hal ini pada tangganya (yakni $F = c$ yang diperluas dengan $0$
        pada $\R_-$ dan $1$ pada $\intco1\infty$): diperolehlah
        \emph{\omterm{def:b3:measure:measure}{ukuran} Cantor} $\mu$. Tunjukkan $\mu(\R) = 1$, bahwa
        setiap celah himpunan Cantor berukuran-$\mu$ nol (sebab $c$
        konstan di sana), lalu simpulkan
        \[
        \mu(C) = 1, \qquad \lambda(C) = 0 :
        \]
        jadi $\mu$ dan $\lambda$ hidup pada pembawa yang saling lepas (yakni $C$ dan
        komplemennya). Dua \omterm{def:b3:measure:measure}{ukuran} dalam kedudukan itu
        disebut \emph{saling singular}, ditulis $\mu \perp
        \lambda$.
  \item Tunjukkan bahwa $\mu$ tak beratom: $\mu(\{x\}) = 0$ untuk
        setiap $x$ \emph{(lewat \omterm{def:b3:topology:continuity}{kekontinuan} $c$)}. Jadi ada \omterm{def:b3:measure:measure}{ukuran}
        peluang tanpa atom yang dibawa sebuah \omterm{def:b3:topology:compact}{kompak} berukuran nol Lebesgue:
        bandingkan dengan \omterm{def:b3:measure:measure}{ukuran} yang sudah kita lihat sejauh ini.
  \item (Pelemparan koin yang menyamar) Untuk sebuah kata $(\varepsilon_1,
        \dots, \varepsilon_m) \in \{0,1\}^m$, misalkan
        $C_{\varepsilon}$ himpunan $x \in C$ yang angka ternernya
        memenuhi $b_i(x) = \varepsilon_i$ untuk $i \leq m$
        (yakni salah satu dari $2^m$ kepingan Cantor berkedalaman $m$). Tunjukkan
        $\mu(C_\varepsilon) = 2^{-m}$ \emph{(sebab tangganya
        mendaki $2^{-m}$ sepanjang kepingan itu: pakailah Bagian I,
        pertanyaan 3)}. Jadi \omterm{def:b3:measure:measure}{ukuran} Cantor merupakan hukum sebuah
        barisan tak berhingga lemparan koin adil yang dibaca dalam terner
        --- dan \cref{ch:b3:probability} akan mengeksakkannya.
  \item Buktikan keserupaan dirinya: bahwa untuk setiap $A$ yang Borel,
        \[
        \mu(A) = \tfrac12\,\mu(3A) +
        \tfrac12\,\mu(3A - 2) ,
        \]
        dengan $3A - 2 = \{3x - 2 : x \in A\}$ \emph{(periksalah
        pada selang pembangkitnya $\intoc ab$ lewat relasi
        serupa-diri $c$, lalu panggillah ketunggalannya,
        \cref{thm:b3:measure:uniqueness})}.
  \item Tunjukkan bahwa pencerminan $s(x) = 1 - x$ mempertahankan
        $\mu$: yakni $\mu(s(A)) = \mu(A)$ \emph{(lewat $c(1 - x) = 1 -
        c(x)$, yang menyusul dari kesimetrian
        konstruksinya --- buktikanlah)}.
  \item Hitunglah dua momen pertama $\mu$, yakni dari sebuah
        titik acak $X$ berhukum $\mu$ (integralnya boleh
        digarap sebagai limit jumlah atas kepingan berkedalaman $m$,
        sambil mendahului \cref{ch:b3:lebesgue}): kesimetriannya memberikan
        $\int x\,\dd\mu = \frac12$, dan keserupaan dirinya memberikan
        \[
        \int x^2\,\dd\mu = \frac38,
        \qquad\text{sehingga}\qquad
        \operatorname{Var}(X) = \frac18 .
        \]
        Bandingkan dengan hukum seragam pada $\intcc01$ (yang beragam
        $\frac1{12}$): jadi massa Cantor, yang terdorong ke tepinya,
        menyebar \emph{lebih} lebar.
  \item Tunjukkan bahwa \emph{pendukung} (secara \omterm{def:b3:topology:topology}{topologi}) dari $\mu$
        --- yakni himpunan tertutup terkecil yang berukuran penuh --- tepat
        berupa $C$.
  \item (Sintesis) Tangga $c$ bersifat \omterm{def:b3:topology:continuity}{kontinu} dan
        tak turun namun menggagalkan teorema fundamental
        kalkulus (Bagian I); sedangkan \omterm{def:b3:measure:measure}{ukuran} $\mu_c$ merupakan
        peluang, tanpa atom, dan singular terhadap
        $\lambda$. Jelaskan dalam satu paragraf pendek bagaimana keduanya
        merupakan dua wajah satu gejala, lalu nyatakan
        pelajaran umumnya: bahwa fungsi tak turun bersesuaian dengan \omterm{def:b3:measure:measure}{ukuran}
        ($F \leftrightarrow \mu_F$), keterdiferensialan hampir di mana-mana
        bersesuaian dengan ``bagian \omterm{def:b3:topology:continuity}{kontinu} mutlaknya'', dan
        $c$ merupakan saksi baku bahwa $F$ yang \omterm{def:b3:topology:continuity}{kontinu} dapat
        \emph{tak} membawa bagian \omterm{def:b3:topology:continuity}{kontinu} mutlak sama sekali.
  \item (Modulus \omterm{def:b3:topology:continuity}{kekontinuan} yang tajam) Misalkan $s =
        \frac{\ln 2}{\ln 3}$. Tunjukkan bahwa $c$ bersifat Hölder
        \omterm{def:b3:topology:continuity}{kontinu} dengan eksponen $s$:
        \[
        \abs{c(x) - c(y)} \leq 4\,\abs{x - y}^{s}
        \qquad (x, y \in \intcc01),
        \]
        dan bahwa tak ada eksponen $t > s$ yang berlaku, bahkan secara lokal.
        Lalu simpulkan bentuk teori \omterm{def:b3:measure:measure}{ukurannya}: bahwa untuk setiap $x$ dan
        setiap $r \in \intoc01$,
        \[
        \mu\bigl(\intcc{x - r}{x + r}\bigr) \leq 8\,r^{s} .
        \]
        \emph{(Bandingkan kisi triadik berkedalaman $m$ dengan skala
        $\abs{x - y}$; dan pertanyaan 18 memberikan pendakian sepanjang
        setiap kepingannya. Eksponen $s$ merupakan dimensi \omterm{def:b3:topology:hausdorff}{Hausdorff}
        $C$, sebagaimana akan dikatakan kuliah berikutnya.)}
  \item (Keserupaan diri mencirikan $\mu$) Buktikan
        konvers pertanyaan 19: jika $\nu$ sebuah \omterm{def:b3:measure:measure}{ukuran} peluang
        pada $\mathcal B(\R)$ yang dibawa $\intcc01$ dan
        memenuhi
        \[
        \nu(A) = \tfrac12\,\nu(3A) + \tfrac12\,\nu(3A - 2)
        \qquad (A \in \mathcal B(\R)),
        \]
        maka $\nu = \mu$. \emph{(Iterasikan relasinya $m$
        kali untuk menyebarkan $\nu$ atas $2^m$ kepingan Cantor
        berkedalaman $m$, taksirlah $\nu(\intoc ab)$ terhadap cacahan
        kepingan di dalam $\intoc ab$, lalu biarkan $m \to \infty$;
        dan selesaikan dengan \cref{thm:b3:measure:uniqueness}.)}

\end{enumerate}
\end{problem}
